\documentclass[10pt,a4paper]{amsart}
\usepackage[margin=19mm]{geometry}
\usepackage{amsmath,amssymb,mathtools}
\usepackage[scr=rsfs,cal=euler]{mathalfa}
\usepackage{xfrac}
\usepackage[unicode,hidelinks,bookmarks=false]{hyperref}
\newcommand{\ie}{i.e.\ }
\newcommand{\cf}{cf.\ }
\newcommand{\resp}{respectively\ }
\newcommand{\Subsection}{\subsection}
\providecommand{\nobreakdash}{\nobreak\hbox{-}}
\setlength{\emergencystretch}{2em}
\multlinegap0pt
\usepackage{bm}
\usepackage{bbm}
\usepackage{dsfont}
\usepackage{xspace}
\usepackage{cancel}
\usepackage{mathtools}
\usepackage{amsthm}
\makeatletter
\@ifundefined{theorem}{\newtheorem{theorem}{Theorem}}{}
\@ifundefined{proposition}{\newtheorem{proposition}[theorem]{Proposition}}{}
\makeatother
\def \text{\mbox}
\renewcommand{\mod}{\mathop{\rm mod}\nolimits}%
\title[Grothendieck groups of repetitive cluster categories]{Grothendieck groups of repetitive cluster categories}
\author{Huimin Chang, Dave Murphy, Panyue Zhou}
\date{Source: arXiv:2501.11021v2 (CC BY 4.0)}
\begin{document}
\maketitle

\noindent\emph{Source and licence.} Grothendieck groups of repetitive cluster categories. Version arXiv:2501.11021v2 (CC BY 4.0), distributed under the Creative Commons Attribution 4.0 International licence, as recorded in the arXiv licence metadata for this exact version and shown on the abstract page https://arxiv.org/abs/2501.11021v2. Full rights notice: licenses/arxiv-2501.11021-CC-BY-4.0.txt.\par\smallskip

\noindent\textit{Editorial scope.} Counted unit: the Theorem whose source label is \texttt{Thm: n odd, p odd} in the article (source0.tex lines 801--821, source label \texttt{Thm: n odd, p odd}), delivered together with the article's own complete original proof of it (source0.tex lines 823--852, one \texttt{proof} environment of 30 lines). No mathematics is written, repaired or re-derived: statement and proof are the article's own text line for line. Required context retained uncounted: Prop: n odd, epi to K0 (lines 740--747). Every definition, display and label of those lines is retained unchanged. Omitted from this package and not redistributed: 1--739, 748--800, 822--822, 853--1303. Nothing inside a retained excerpt is omitted or altered: every retained body is delivered exactly as the article's lines are written, so no omission receipt of any kind accompanies this package. Citations: the retained excerpts carry no citation command at all, so no bibliography entry of the article is transcribed and no reference is redistributed. Reference adaptation: none was needed and none was made; every reference inside the delivered unit resolves inside it, so no unresolved reference or citation is produced. Printed slips kept literal: no printed word, symbol, hypothesis or number of the counted statement or of its proof was altered. Layout and class: the article's own class is \texttt{amsart} with packages including amsmath, amssymb, amsthm, latexsym, xypic, amsfonts, mathrsfs, graphicx, color, xcolor, tikz, graphics, hyperref, anysize; the wrapper is the lane's pinned \texttt{amsart} editorial wrapper at 10pt on A4, which restates the theorem-like environments the retained text needs and adds the article's own macro definitions that the retained text needs (\texttt{\textbackslash mod}, \texttt{\textbackslash text}), with extra wrapper packages (bm, bbm, dsfont, xspace, cancel, mathtools). The delivered numbering is the unit's own. Assets: no retained span contains a figure environment, a graphics-inclusion command, a PDF-page inclusion, a table or a TikZ picture; no image file is copied into this package, the package declares no asset dependency and no image file is redistributed with it. Licence: version arXiv:2501.11021v2 of this article is distributed under the Creative Commons Attribution 4.0 International licence, as recorded in the arXiv licence metadata for this exact version and shown on the abstract page https://arxiv.org/abs/2501.11021v2. A full rights notice is delivered at licenses/arxiv-2501.11021-CC-BY-4.0.txt. Characters: every retained excerpt is pure ASCII, so the delivered engine, XeLaTeX with its default Latin Modern text font, reports no missing character.
\begin{proposition}\label{Prop: n odd, epi to K0}
	Let $n$ be odd, suppose that $0 < p \leq \frac{n+1}{2}$ and $n+1 \equiv k \mod (p)$ with $0 < k \leq p$.
	Then there exists an epimorphism
	\[
	\langle [S_1], [S_2],\ldots, [S_m] \rangle \twoheadrightarrow K_0(C_{n,p}^A),
	\]
	where $m$ is the greatest common divisor of $p$ and $k$.
\end{proposition}

\begin{theorem}\label{Thm: n odd, p odd}
	Let $n,p,m$ as in Proposition \ref{Prop: n odd, epi to K0}, and further suppose that $p \not\equiv \pm 1 \mod (n+1)$ is odd, and that $\frac{p}{m}$ is even.
	Then,
\[ K_0(C_{n,p}^A) =\left\{
\begin{array}{ll}
\langle [S_1], [S_2], \ldots, [S_m] \mid (-1)^t [S_c] - \sum_{j=1}^{b} [S_j] = 0 \rangle, & \text{if\;} $a$ \text{\;is\;even},\\
		\langle [S_1], [S_2], \ldots, [S_m] \mid (-1)^t [S_c] - \sum_{j=b+1}^{m} [S_j] = 0 \rangle, & \text{if\;} $a$ \text{\;is\;odd},
\end{array}
\right.
\]	
	where $n+1-p = tm +c$, and $am+b=n$.
	Further, if $\frac{p}{m}$ is odd.
	Then,
\[ K_0(C_{n,p}^A) =\left\{
\begin{array}{ll}
\langle [S_1], [S_2], \ldots, [S_m] \mid 2[S_j] = 0, (-1)^t [S_c] - \sum_{j=1}^{b} [S_j] = 0\rangle, & \text{if\;} $a$ \text{\;is\;even},\\
		\langle [S_1], [S_2], \ldots, [S_m] \mid 2[S_j] = 0, (-1)^t [S_c] - \sum_{j=b+1}^{m} [S_j] = 0 \rangle, & \text{if\;} $a$ \text{\;is\;odd}.
\end{array}
\right.
\]	
\end{theorem}

\begin{proof}
	As $p$ is odd, then our relations are $[S_j] + [S_{j+m}]$ by Proposition \ref{Prop: n odd, epi to K0}, and so $[S_j] = (-1)^s[S_{j+sm}]$ in $K_0(C_{n,p}^A)$.
	
	We have two conditions to consider, whether $m \mid p$ is odd or even, and if $a$ is even.
	\begin{itemize}
		\item 	If $\frac{p}{m}$ is even, then there exists an even number of elements in the equivalence class of $[S_i]$ in $\{[S_1],[S_2], \ldots, [S_p]\}$ for $0 < i \leq m$, and so $[S_i] = (-1)^s[S_{i+sm}]$, where $s$ is even if $sm=p$.
		Therefore we do not induce any more relations on $[S_i]$ than those we have already considered, as $[S_i] = [S_i]$.
		
		However, if $\frac{p}{m}$ is odd, then there exists an odd number of elements in the equivalence class of $[S_i]$ in $\{[S_1],[S_2], \ldots, [S_p]\}$ for $0 < i \leq m$, and so $[S_i] = (-1)^s[S_{i+sm}]$, where $s$ is odd if $sm=p$.
		Therefore we find the new relation $[S_i] = -[S_i]$ for all $0 < i \leq m$, and so $2[S_i]=0$.\\
		
		\item Consider the relation $[S_{n+1-p}] - \sum_{j=1}^n [S_j]$.
		Then $[S_{n+1-p}] = (-1)^t[S_c]$ for $tm + c = n+1-p$, and $\sum_{j=1}^n [S_j] = \sum_{j=1}^r [S_j]$, where $2tm + r =n$ for $0 \leq r < 2m$.
		This is due to
		\[
		\sum_{j=r+1}^{r+1+2m} [S_j] = \sum_{j=r+1}^{r+1+m} [S_j] + [S_{j+m}] = \sum_{j=r+1}^{r+1+m} [S_j] - [S_j] = 0.
		\]
		
		If $a$ is even, then $a=2t$ and $r=b < m$, and so we have the relation $[S_{n+1-p}] - \sum_{j=1}^n [S_j] = [S_{n+1-p}] - \sum_{j=1}^b [S_j] = (-1)^t[S_c] - \sum_{j=1}^b [S_j]$.
		Hence $(-1)^t[S_c] - \sum_{j=1}^b [S_j] =0$ in $K_0(C_{n,p}^A)$ if $a$ is even.
		
		If $a$ is odd, then $a=2t+1$, and $r=b+m$, so we have the relation
		\[
		[S_{n+1-p}] - \sum_{j=1}^n [S_j] = (-1)^t[S_c] - \sum_{j=1}^m [S_j] + \sum_{j=1}^b [S_j] = (-1)^t [S_c] - \sum_{j=b+1}^m [S_j].
		\]
		Hence $(-1)^t [S_c] - \sum_{j=b+1}^m [S_j] = 0$ in $K_0(C_{n,p}^A)$ if $a$ is odd.
	\end{itemize}
	
	Our claim follows by consideration of all combinations of the above two cases.
\end{proof}

\end{document}
